

```latex
\documentclass[11pt]{article}

\usepackage{amsmath,amssymb,amsthm}
\usepackage[margin=1in]{geometry}

\theoremstyle{plain}
\newtheorem{theorem}{Theorem}[section]
\newtheorem{lemma}[theorem]{Lemma}
\newtheorem{corollary}[theorem]{Corollary}
\newtheorem{proposition}[theorem]{Proposition}

\theoremstyle{definition}
\newtheorem{definition}[theorem]{Definition}
\newtheorem{remark}[theorem]{Remark}

\numberwithin{equation}{section}

\title{Prime summands of Dehn surgeries on L-space knots}
\author{}
\date{}

\begin{document}

\maketitle

\begin{abstract}
We prove that every closed 3-manifold obtained by Dehn surgery on an L-space knot has at most two prime summands.  In fact, we prove a stronger classical statement: every closed 3-manifold obtained by Dehn surgery on any knot in $S^3$ has at most two prime summands.  The proof uses only standard results on reducible Dehn surgery, namely Scharlemann's theorem on the integrality of reducible surgery slopes and the Gordon--Luecke cabling theorem together with the standard cable-surgery formula.
\end{abstract}

\section{Conventions and preliminaries}

We use the standard notation $S^3_{a/b}(K)$ for the closed 3-manifold obtained by $a/b$-Dehn surgery on a knot $K\subset S^3$, where $a/b\in \mathbb Q\cup\{\infty\}$ is written with $\gcd(a,b)=1$ and $b\ge 0$, and $\infty=1/0$ denotes the meridional filling.  Integer surgery $n$ means slope $n/1$.

By an \emph{L-space knot} we mean a knot $K\subset S^3$ such that some positive Dehn surgery on $K$ is an L-space \cite{OSz}.  Many authors require L-space knots to be nontrivial; our argument also covers the unknot.

A closed orientable 3-manifold is \emph{irreducible} if every embedded $2$-sphere bounds a $3$-ball.  It is \emph{prime} if whenever it is written as a connected sum $M_1\# M_2$, one of the summands is homeomorphic to $S^3$.  By the Kneser--Milnor prime decomposition theorem, every closed orientable 3-manifold has a finite connected-sum decomposition into prime manifolds, unique up to order and omission of $S^3$ factors \cite{Hempel}.  In this note, ``prime summands'' means prime factors other than $S^3$.

We record a few elementary facts.

\begin{lemma}\label{lem:basic}
\begin{enumerate}
\item Every closed orientable irreducible 3-manifold is prime.
\item Lens spaces are prime.  The manifold $S^2\times S^1$ is also prime.
\item If $U\subset S^3$ is the unknot, then $S^3_{a/b}(U)$ is a lens space when $a\ne 0$, and is $S^2\times S^1$ when $a=0$.  In particular, surgery on the unknot has at most one prime summand.
\end{enumerate}
\end{lemma}

\begin{proof}
The first assertion is immediate: if an irreducible manifold were a nontrivial connected sum, the sphere along which the connected sum is formed would be an essential sphere.

Lens spaces are irreducible; for example, they have universal cover $S^3$, so $\pi_2=0$, and by the sphere theorem they contain no essential sphere.  Hence they are prime.  The fact that $S^2\times S^1$ is prime is standard; see for instance \cite{Hempel}.

Finally, Dehn filling a solid torus along a torus of slope $a/b$ gives the lens space $L(a,b)$ if $a\ne 0$, and gives $S^2\times S^1$ if $a=0$ \cite{Rolfsen}.  Lens spaces and $S^2\times S^1$ are prime, so surgery on the unknot has at most one prime summand.
\end{proof}

\section{Reducible surgeries on knots in $S^3$}

We now recall the two deep classical inputs.

\begin{theorem}[Scharlemann {\cite{Sch}}]\label{thm:Scharlemann}
Let $J\subset S^3$ be a nontrivial knot.  If $S^3_{a/b}(J)$ is reducible, then $|b|=1$.  Equivalently, every nonintegral Dehn surgery on a nontrivial knot in $S^3$ is irreducible.
\end{theorem}

\begin{theorem}[Gordon--Luecke {\cite{GL}}]\label{thm:GL}
Let $K\subset S^3$ be a nontrivial knot.  If $S^3_r(K)$ is reducible, then $K$ is a nontrivial cable knot.  More precisely, $K$ is the $(p,q)$-cable of a companion knot $C\subset S^3$, with $|p|\ge 2$ and $\gcd(p,q)=1$, where $|p|$ is the winding number of the cable.  The reducible surgery slope is the integer cabling slope
\[
r=pq,
\]
and there is a homeomorphism
\begin{equation}\label{eq:cable-formula}
S^3_{pq}(K)\cong L\#\, S^3_{q/p}(C),
\end{equation}
where $L$ is a lens space.
\end{theorem}

\begin{remark}
The connected-sum decomposition \eqref{eq:cable-formula} is the standard cable-surgery formula; see also \cite{Rolfsen}.  The precise orientation and the precise lens space $L$ depend on conventions, but the only properties used here are that $L$ is a lens space and that the induced surgery slope on the companion is $q/p$, whose denominator is the winding number $|p|$.
\end{remark}

\section{The main theorem}

\begin{theorem}\label{thm:all-knots}
Let $K\subset S^3$ be any knot, and let $r\in \mathbb Q\cup\{\infty\}$ be any surgery slope.  Then $S^3_r(K)$ has at most two prime summands.  Moreover, if $K$ is nontrivial and $S^3_r(K)$ is reducible, then it is a connected sum of a lens space and a prime manifold.
\end{theorem}

\begin{proof}
First suppose $K$ is the unknot.  By Lemma \ref{lem:basic}(3), $S^3_r(K)$ is either a lens space or $S^2\times S^1$, hence has at most one prime summand.

Now suppose $K$ is nontrivial, and put $M=S^3_r(K)$.  If $M$ is irreducible, then by Lemma \ref{lem:basic}(1) it is prime, so it has at most one prime summand.

Assume therefore that $M$ is reducible.  By Theorem \ref{thm:GL}, $K$ is a nontrivial $(p,q)$-cable of a companion knot $C$, with $|p|\ge 2$ and $\gcd(p,q)=1$, the surgery slope is $r=pq$, and
\[
M\cong L\#\, S^3_{q/p}(C),
\]
where $L$ is a lens space.  By Lemma \ref{lem:basic}(2), $L$ is prime.

We claim that $Y:=S^3_{q/p}(C)$ is prime.  If $C$ is the unknot, then $Y$ is a lens space by Lemma \ref{lem:basic}(3), hence prime.  If $C$ is nontrivial, then $q/p$ is not an integer, because $\gcd(p,q)=1$ and $|p|\ge 2$.  Therefore Theorem \ref{thm:Scharlemann} implies that $Y$ is irreducible, and Lemma \ref{lem:basic}(1) implies that $Y$ is prime.

Thus $M$ is a connected sum of two prime manifolds, up to possible $S^3$ factors which are omitted in the prime decomposition.  Hence $M$ has at most two prime summands.
\end{proof}

\begin{corollary}\label{cor:L-space-knot}
Let $K\subset S^3$ be an L-space knot.  Then every closed 3-manifold obtained by Dehn surgery on $K$ has at most two prime summands.
\end{corollary}

\begin{proof}
By definition, an L-space knot is a knot in $S^3$.  The result is therefore immediate from Theorem \ref{thm:all-knots}.
\end{proof}

\begin{remark}
The proof shows that the L-space hypothesis is not used.  The statement is a purely topological consequence of the classification of reducible surgeries on knots in $S^3$.  In particular, the only way a surgery on a nontrivial knot in $S^3$ can have two prime summands is when the knot is a cable knot and the surgery is the cabling surgery; one summand is a lens space, and the other is an irreducible rational surgery on the companion knot.
\end{remark}

\begin{thebibliography}{9}

\bibitem{Hempel}
J.~Hempel, \emph{3-Manifolds}, Annals of Mathematics Studies, No.~86, Princeton University Press, Princeton, NJ, 1976.

\bibitem{Rolfsen}
D.~Rolfsen, \emph{Knots and Links}, Mathematics Lecture Series, No.~7, Publish or Perish, Berkeley, CA, 1976.

\bibitem{Sch}
M.~Scharlemann, Producing reducible 3-manifolds by surgery on a knot, J. Differential Geom. \textbf{22} (1985), no.~1, 49--73.

\bibitem{GL}
C.~McA. Gordon and J.~Luecke, Dehn surgery on knots, Ann. of Math. (2) \textbf{129} (1989), no.~2, 199--248.

\bibitem{OSz}
P.~Ozsv\'ath and Z.~Szab\'o, Absolutely graded Floer homologies and intersection forms for four-manifolds with boundary, Adv. Math. \textbf{173} (2003), no.~2, 179--261.

\end{thebibliography}

\end{document}
```